\section{Stress, anxiety and the prefrontal executive system}\label{sec:stress}
\subsection{Mechanistic neuroscience}
The most developed account is due to Arnsten. Her review of stress signalling describes the prefrontal cortex (PFC) as the most evolved brain region and also the one most sensitive to the detrimental effects of stress exposure \citep{arnsten2009stress}. The proposed mechanism is an inverted-U in catecholamine signalling: moderate norepinephrine and dopamine levels support prefrontal network firing and working memory, while the high levels of uncontrollable stress weaken them and strengthen amygdala and striatal habit circuitry. A later review summarises the result as stress rapidly flipping the brain from reflective to reflexive control of behaviour \citep{arnsten2015stress}. Acute uncontrollable stress impairs PFC-dependent working memory in monkeys and rodents \citep{arnsten1998noise}. In humans, a meta-analysis finds that acute stress impairs working memory and cognitive flexibility, with more nuanced effects on inhibition and many moderators \citep{shields2016stress}, and imaging suggests a reallocation of resources toward salience networks rather than a switch-off \citep{hermans2014networks}.

\paragraph{Where the ``off-line'' language comes from} The phrase traces to a hedged line in a 1998 monkey study (``Stress may take the PFC `off-line'''; \citealt{arnsten1998noise}), and its author still uses it in quotation marks for \emph{uncontrollable} stress \citep{arnsten2020guanfacine}. Brewer and Roy apply it to worry and cite Arnsten's 2009 and 2015 reviews, which are not studies of worry \citep{brewer2021habit}; their app-based trial ($n=65$, treatment-as-usual control) did not test a prefrontal mechanism \citep{roy2021app}. We therefore read ``offline'' as a popularisation of \emph{impaired, in a graded and dose-dependent way}, mainly under uncontrollable stress.

\subsection{Does stress abolish reappraisal in humans?}
Evidence for this specific link is contested. The founding study found that participants stressed before learning showed no fear reduction after regulation training \citep{raio2013stress}, and a review summarises that acute stress may impair the persistent inhibition of fear, presumably by altering prefrontal function \citep{raio2015influence}. Five later studies ($n=54$--$98$) found no group-level impairment \citep{shermohammed2017stress,sandner2021withstands,langer2021delayed,langer2023rapid,hamza2024acceptance}, though two reported sex- or cortisol-linked effects. Their laboratory stressors were moderate, so the null results may not transfer to clinical-level arousal; that gap is itself an open question.

\subsection{Cognitive psychology of anxiety}
Attentional Control Theory proposes that anxiety impairs the inhibition and shifting functions of executive control, so that anxious individuals are more captured by threat and worse at redirecting attention \citep{eysenck2007act}. The theory predicts that \emph{efficiency} (effort and time) suffers more than \emph{effectiveness} (accuracy), because anxious people compensate with extra effort. Meta-analytic support exists at group level \citep[$g=-0.58$;][]{shi2019attentional} and for working memory \citep[$g=-0.334$;][]{moran2016working}, but with qualifiers: a meta-analysis of 32 studies ($N=13{,}084$) found generalised anxiety disorder (GAD) associated with poorer cognitive flexibility and working memory but not inhibitory control, and with no association with any executive domain in older adults \citep{nguyen2025executive}. Working memory was unrelated to social anxiety in one study \citep{harrison2025executive}, and one study found \emph{enhanced} control over negative distractors in GAD \citep{hallion2019enhanced}.
